\section{Method}

\subsection{Task, observations, and actions}

Push-T fixes the goal at world position $(-0.156,-0.100)$ m with yaw $300^\circ$ and samples the
T block across position and orientation~\cite{maniskill_pusht}. Success requires the rendered
intersection between the block and goal to reach at least 0.90. The primary metric,
\successonce, records whether this threshold is reached at any point; \successend records whether
it remains satisfied at the time limit.

The policy receives RGB images and robot proprioception, but no privileged block or goal pose. At
each decision it conditions on observation horizon $T_o=2$. RGB frames are scaled to $[0,1]$ and
processed by a five-layer convolutional encoder. A $128\!\times\!128$ image becomes a
$128\!\times\!8\!\times\!8$ feature map, which is flattened and projected to 256 dimensions.
Retaining this spatial map is important because global pooling discards absolute block and goal
location. The visual feature is concatenated with robot state at each observation step.

The demonstrations and evaluation both use ManiSkill's \texttt{pd\_ee\_delta\_pose} controller.
Actions are normalized to $[-1,1]$. We execute one action and immediately replan
($T_a=1$), matching the tuned Push-T control pattern and limiting open-loop error during contact.

\subsection{Conditional action diffusion}

The policy predicts noise for a $T_p=16$ action sequence with a conditional one-dimensional
U-Net. It uses channel widths $[64,128,256]$, GroupNorm, Mish activations, skip connections, and
FiLM-style conditioning on visual-proprioceptive history and a 64-dimensional sinusoidal
diffusion-step embedding. Training samples diffusion step $k$, adds noise $\epsilon$ to a
demonstrated action sequence, and minimizes
\begin{equation}
    \mathcal{L}(\theta)=
    \mathbb{E}_{a,\epsilon,k}\!
    \left[\left\|\epsilon-\epsilon_\theta(a_k,k,o_{t-1:t})\right\|_2^2\right].
\end{equation}
We use 100 denoising steps and a squared-cosine DDPM schedule~\cite{nichol2021improved}.

\subsection{Aligned RGB demonstration replay}
\label{sec:aligned-replay}

The downloaded RL demonstrations contain simulator states rather than the required RGB sequence.
We replay the first 100 successful \texttt{pd\_ee\_delta\_pose} trajectories with
\texttt{physx\_cuda}, RGB observations, and the original 1,024 parallel environments. The
validated dataset contains 9,865 transitions and 9,765 training windows.

A trajectory with $T$ actions contains $T+1$ environment states. After action $a_t$, the
observation paired with the next action must begin at $s_{t+1}$. The audited upstream replay path
restored $s_t$ instead; our fail-closed wrapper changes this operation to
\texttt{env\_states[t + 1]} in process and rejects source code that does not match the audited
pattern. Before training, the validator checks source success labels, controller, collection and
replay parallelism, output observation mode, trajectory count, and exactly $T+1$ RGB observations
for $T$ actions in every episode.

\subsection{Training and checkpoint protocol}

We train one model with seed 1 using AdamW, batch size 128, learning rate $10^{-4}$, 500 linear
warmup steps, and cosine decay. Demonstration tensors remain on CPU and minibatches are transferred
to the GPU. EMA weights with power 0.75 are used for rollout evaluation. Training runs for 50,000
updates; loss is logged every 100 updates, and a checkpoint is saved every 5,000.

Every 5,000 updates, 20 fresh episodes provide a model-selection diagnostic. A checkpoint replaces
\texttt{best\_success\_once.pt} only when \successonce strictly improves. Interrupted sessions
resume the complete training state, not just model weights. The run begins from a 10k Colab
checkpoint, continues to 15k, and resumes from the validated 15k checkpoint to 50k on Runpod.

\subsection{Fixed-checkpoint evaluation}

Final evaluation loads the single selected EMA checkpoint and runs 100 fresh episodes for each
environment seed in $\{0,1,2\}$. All seeds use the same training checkpoint; they are not three
independently trained policies. The evaluator reconfigures every environment, ignores early task
termination until the 150-step limit, and records synchronized complete episodes. We report each
seed, the arithmetic mean and sample standard deviation of the three success rates, and pooled
successes out of 300. Auxiliary metrics include return, maximum overlap, final overlap, and
\successend.
